\documentclass[11pt]{article}
\usepackage[a4paper,margin=24mm]{geometry}
\usepackage{amsmath,amssymb,amsthm}
\usepackage[hidelinks]{hyperref}
\newtheorem{theorem}{Theorem}
\newtheorem{lemma}{Lemma}
\begin{document}
\title{A counterexample to the explicit period condition\newline in Conjecture 00000000231}
\author{Chengzhe Gao}\date{}\maketitle

\paragraph{Source and scope.}
Both language versions define the exceptional primes by the condition
that the Fibonacci period does not divide $p^2$, and assert that none
is below $10^{17}$. Write $F_0=0$, $F_1=1$ and
$F_{n+2}=F_n+F_{n+1}$, and let $\pi(m)$ be the least positive period
of the entire Fibonacci sequence modulo $m$.
Under the explicit condition $\pi(p)\nmid p^2$, the prime $p=7$
is a counterexample. The same prime works if the unspecified period
is instead read as $\pi(p^2)$.

\begin{theorem}
The prime $7<10^{17}$ satisfies
\[
 \pi(7)=16,\qquad \pi(49)=112,\qquad
 16\nmid49,\qquad112\nmid49.
\]
Thus the source's universal nonexistence clause is false under either
of the two modulus readings above.
\end{theorem}

\begin{lemma}[A finite certificate for an infinite sequence]
Let $R_m(n)=(F_n,F_{n+1})\pmod m$. If $R_m(k)=R_m(0)$ and $k>0$,
then $F_{n+k}\equiv F_n\pmod m$ for every $n\geq0$.
Conversely, any positive period $k$ has $R_m(k)=R_m(0)$.
In particular, if there is no such return at $1\leq j<k$, then
$\pi(m)=k$.
\end{lemma}
\begin{proof}
The residue-pair recurrence is the deterministic map
$(a,b)\mapsto (b,a+b\pmod m)$. Equality of the pairs at $k$ and $0$
therefore gives equality at $n+k$ and $n$ by induction on $n$.
The converse uses the period equalities at $n=0$ and $n=1$.
Any smaller positive period would consequently give a smaller return.
\end{proof}

\begin{proof}[Proof of the theorem]
Modulo $7$, the residues $F_0,\ldots,F_{17}$ are
\[
0,1,1,2,3,5,1,6,0,6,6,5,4,2,6,1,0,1.
\]
Among the pairs with starting indices $1,\ldots,16$, the first return
to $(0,1)$ is at $16$. The lemma gives $\pi(7)=16$.

For completeness, $\pi(49)=112$ can be checked with the same exact
finite recurrence. A short certificate, for either modulus, is:
\begin{verbatim}
a, b = 0, 1
for j in range(1, k + 1):
    a, b = b, (a + b) % m
    assert ((a, b) == (0, 1)) == (j == k)
\end{verbatim}
For $(m,k)=(7,16)$ and $(49,112)$ every assertion holds.
The Lean project verifies every earlier pair and the terminal pair
by kernel reduction of the natural-number Fibonacci recurrence,
then applies the lemma. Thus this certificate establishes the
least period of the entire sequence, with no assumption about
unexamined indices. The supplied Python script independently
recomputes both residue orbits as a supplementary check.

Finally $7$ is prime, $7<10^{17}$, and neither $16$ nor $112$ divides
$49$, proving both stated counterexamples.
\end{proof}

\paragraph{Relation to the conventional problem.}
The source's divisibility condition is different from the
conventional exceptional-prime condition $\pi(p^2)=\pi(p)$;
see Elsenhans and Jahnel~\cite{EJ}, Section~2.2.
Indeed, the calculations above give $112\ne16$, so $7$ is
\emph{not} an exceptional prime in that conventional sense.
This submission only disproves the explicit bilingual statement.
It makes no claim about the existence or infinitude of conventional
Wall--Sun--Sun primes or the validity of searches for them.
The later discriminant and infinitude clauses are unnecessary to
the disproof of the first conjunct.

\paragraph{Formalization and reproduction.}
The Lean 4.19.0 project uses only its bundled standard library.
It defines the actual Fibonacci sequence on all natural indices,
proves its initial values, recurrence and uniqueness, and defines
a period by a universal equality over every index.
The generic least-period certificate proves the bridge from
finite residue pairs to this unbounded definition.
Actual natural-number primality, divisibility, leastness and the
bound $10^{17}$ appear in the final counterexample theorems.
Both modulus readings and the inequality of the conventional
periods are formally proved. Kernel-checked \texttt{decide}
certifies only the finite arithmetic; no oracle computation is used.
No \texttt{sorry}, \texttt{admit}, \texttt{native\_decide} or added
axiom occurs. From \texttt{lean/}, run
\texttt{lake build} and
\texttt{lake env lean -DwarningAsError=true Main.lean}.
Build logs, printed axiom dependencies and a solo adversarial
review accompany this submission. No independent-review claim is made.

\begin{thebibliography}{1}
\bibitem{EJ} A.-S. Elsenhans and J. Jahnel,
\emph{The Fibonacci sequence modulo $p^2$---An investigation by
computer for $p<10^{14}$}, 2010.
\href{https://arxiv.org/abs/1006.0824}{arXiv:1006.0824}.
\end{thebibliography}
\end{document}
